% =========================================================================
% Chapter 3: Materials and Methods
% =========================================================================

\chapter{Materials and Methods}
\label{chap:methods}

\section{Taxon Sampling and Specimen Examination}
\label{sec:methods_sampling}

Field surveys were conducted across 18 geographic localities situated in the department of San Mart{\'i}n and Loreto in northeastern Peru, encompassing elevations from \SI{180}{\metre} to \SI{1150}{\metre} above sea level (\cref{app:specimens}). Animal capture, ethical handling, and voucher preservation strictly conformed to the institutional protocol approved by the Institutional Animal Care and Use Committee (\texttt{IACUC-2023-BIO-0492-A}) and statutory collection permits issued by SERFOR and MAATE.

A total of 124 adult museum voucher specimens were examined from the collections of the American Museum of Natural History (AMNH), National Museum of Natural History (USNM), Museo de Zoolog{\'i}a de la Pontificia Universidad Cat{\'o}lica del Ecuador (QCAZ), and the Museo de Historia Natural de la Universidad Nacional Mayor de San Marcos (MUSM). Complete locality metadata, geographical coordinates (WGS84 datum), museum catalog numbers, and corresponding GenBank accession numbers are compiled in \cref{app:specimens}.

\section{Morphometric Data Acquisition and Morphological Characters}
\label{sec:methods_morphometrics}

Morphological measurements were taken on adult preserved specimens using digital calipers (Mitutoyo 500-196-30) to the nearest \SI{0.01}{\milli\metre} under a stereomicroscope (Leica S9D). Ten standardized external anatomical variables were measured following \textcite{brown2011ranitomeya}:

\begin{enumerate}
  \item \textbf{Snout-Vent Length (\textsc{svl}):} Distance from tip of snout to posterior cloacal margin.
  \item \textbf{Head Width (\textsc{hw}):} Greatest width of head measured at the level of the tympana.
  \item \textbf{Head Length (\textsc{hl}):} Distance from tip of snout to the posterior angle of the jaw articulation.
  \item \textbf{Eye Diameter (\textsc{ed}):} Horizontal distance between anterior and posterior orbital borders.
  \item \textbf{Tympanum Diameter (\textsc{td}):} Maximum horizontal diameter of the tympanic annulus.
  \item \textbf{Interorbital Distance (\textsc{iod}):} Minimum dorsal distance between upper eyelids.
  \item \textbf{Tibia Length (\textsc{tl}):} Distance from outer knee joint to heel articulation.
  \item \textbf{Foot Length (\textsc{fl}):} Distance from proximal margin of inner metatarsal tubercle to tip of Toe IV.
  \item \textbf{Hand Length (\textsc{hanl}):} Distance from proximal border of palmar tubercle to tip of Finger III.
  \item \textbf{Finger III Disc Width (\textsc{fdw}):} Maximum transverse width of the digital expanded disc on Finger III.
\end{enumerate}

\section{Molecular Protocols: DNA Extraction and Chemical Buffers}
\label{sec:methods_molecular}

Genomic DNA was isolated from liver or thigh muscle tissue biopsies preserved in molecular-grade 96\% ethanol. Tissues were extracted using a modified hexadecyltrimethylammonium bromide (\textsc{ctab}) and chloroform-isoamyl alcohol precipitation protocol \parencite{sambrook2001cloning}.

\begin{protocol}[CTAB DNA Extraction Buffer Preparation]
\label{prot:ctab_buffer}
The working extraction lysis buffer was prepared as follows:
\begin{itemize}
  \item \SI{100}{\milli\molar} \ce{Tris-HCl} (\ce{pH 8.0})
  \item \SI{1.4}{\mole\per\litre} Sodium Chloride (\ce{NaCl})
  \item \SI{20}{\milli\molar} Ethylenediaminetetraacetic Acid (\ce{EDTA}, \ce{pH 8.0})
  \item \SI{2.0}{\percent} (w/v) Hexadecyltrimethylammonium bromide (\textsc{ctab})
  \item \SI{0.2}{\percent} (v/v) $\beta$-mercaptoethanol (\ce{C2H6OS}) added freshly prior to tissue digestion
\end{itemize}
Tissue fragments (\SI{15}{\milli\gram}) were incubated in \SI{600}{\micro\litre} of \textsc{ctab} buffer supplemented with \SI{25}{\micro\litre} of Proteinase K (\SI{20}{\milli\gram\per\milli\litre}) at \SI{56}{\celsius} for 4 hours with continuous gentle agitation at \SI{450}{\rpm}. Following two washes in chloroform:isoamyl alcohol (24:1), DNA was precipitated with ice-cold isopropanol at \SI{-20}{\celsius} overnight, washed twice in 70\% ice-cold ethanol, air-dried, and resuspended in \SI{50}{\micro\litre} of sterile \textsc{te} buffer (\SI{10}{\milli\molar} \ce{Tris-HCl}, \SI{1}{\milli\molar} \ce{EDTA}, \ce{pH 8.0}).
\end{protocol}

\section{Polymerase Chain Reaction (PCR) Amplification}
\label{sec:methods_pcr}

Three gene fragments were targeted for sequencing:
\begin{enumerate}
  \item Mitochondrial Cytochrome $c$ Oxidase Subunit I (\gene{COI}, $\sim$\SI{658}{bp}) using universal primers LCO1490 and HCO2198 \parencite{folmer1994primers}.
  \item Mitochondrial 16S ribosomal RNA (\gene{16S rRNA}, $\sim$\SI{550}{bp}) using primers 16Sar-L and 16Sbr-H.
  \item Nuclear Recombination Activating Gene 1 (\gene{RAG1}, $\sim$\SI{820}{bp}) using primers RAG1-F1 and RAG1-R1.
\end{enumerate}

Detailed oligonucleotide primer sequences, binding orientations, and amplicon lengths are compiled in \cref{app:primers}.

\begin{table}[htbp]
  \centering
  \small
  \caption{Standard PCR master mix formulation for a single \SI{25}{\micro\litre} reaction volume.}
  \label{tab:pcr_mix}
  \begin{tabular}{lcc}
    \toprule
    \textbf{Component} & \textbf{Stock Concentration} & \textbf{Final Concentration} \\
    \midrule
    Sterile Nuclease-Free \ce{ddH2O} & -- & to \SI{25.0}{\micro\litre} \\
    Reaction Buffer (\ce{Tris-HCl}, \ce{KCl}) & $10\times$ & $1\times$ \\
    Magnesium Chloride (\ce{MgCl2}) & \SI{50}{\milli\molar} & \SI{1.5}{\milli\molar} \\
    Deoxynucleotide Triphosphates (\ce{dNTPs}) & \SI{10}{\milli\molar} each & \SI{200}{\micro\molar} each \\
    Forward Primer & \SI{10}{\micro\molar} & \SI{0.4}{\micro\molar} \\
    Reverse Primer & \SI{10}{\micro\molar} & \SI{0.4}{\micro\molar} \\
    Taq DNA Polymerase & \SI{5.0}{\unitU\per\micro\litre} & \SI{0.625}{\unitU} \\
    Template Genomic DNA & $\sim$\SI{25}{\nano\gram\per\micro\litre} & $\sim$\SI{50}{\nano\gram} \\
    \bottomrule
  \end{tabular}
\end{table}

Thermal cycling profiles were executed on an automated thermal cycler (Bio-Rad C1000 Touch). For the mitochondrial \gene{COI} Folmer locus: initial denaturation at \SI{94}{\celsius} for \SI{3}{\minute}; followed by 35 cycles of denaturation at \SI{94}{\celsius} for \SI{30}{\second}, primer annealing at \SI{50}{\celsius} for \SI{45}{\second}, and extension at \SI{72}{\celsius} for \SI{60}{\second}; with a final terminal extension at \SI{72}{\celsius} for \SI{7}{\minute}.

PCR amplicons were visualized by gel electrophoresis on \SI{1.5}{\percent} agarose gels cast in $1\times$ \textsc{tae} buffer (\SI{40}{\milli\molar} \ce{Tris}, \SI{20}{\milli\molar} acetic acid, \SI{1}{\milli\molar} \ce{EDTA}) stained with GelRed (Biotium). Bidirectional Sanger sequencing was performed using BigDye Terminator v3.1 chemistry on an ABI 3730xl Genetic Analyzer.

\section{Phylogenetic Inference and Species Delimitation}
\label{sec:methods_phylogenetics}

Raw chromatograms were edited and assembled in \textsc{Geneious Prime}. Multiple sequence alignments were generated using \textsc{muscle} \parencite{edgar2004muscle} under default gap penalties and visually verified. The concatenation alignment comprised \SI{2028}{bp} across the three gene fragments.

Best-fit models of nucleotide substitution for each gene codon partition were determined using PartitionFinder 2 via the Bayesian Information Criterion (\textsc{bic}). Maximum Likelihood (\textsc{ml}) phylogenetic reconstruction was conducted in \textsc{iq-tree} 2 \parencite{nguyen2015iqtree} with 1,000 ultrafast bootstrap replicates (\textsc{ufboot}) and 1,000 replicates of the Shimodaira-Hasegawa-like approximate likelihood ratio test (\textsc{sh-alrt}).

Bayesian Inference (\textsc{bi}) was executed in \textsc{MrBayes} v3.2.7 \parencite{ronquist2003mrbayes}. Two independent MCMC runs each comprising four chains (one cold, three heated at temperature $T = 0.15$) were executed for 10,000,000 generations, sampling parameters every 1,000 generations. Convergence was diagnosed by confirming average standard deviation of split frequencies fell below 0.005 and effective sample sizes (\textsc{ess}) for all parameters exceeded 200 using Tracer v1.7. The initial 25\% of sampled trees were discarded as burn-in.

Species boundaries were evaluated using multi-rate Poisson Tree Processes (\textsc{mptp}) \parencite{zhang2013ptp} implemented on the ML topology, and through Bayesian species delimitation using \textsc{bpp} v4.4 \parencite{rannala2003bpp} evaluating species model probabilities under the A10 species delimitation algorithm.
